% The Figure structure intentionally includes the non-floating caption.
\tagstructbegin{tag=Figure,alttext={开了的 Query 为一一，三条 Key 匹配分数依次为一、二、三；Softmax 权重为零点零九零、零点二四五、零点六六五，Value 加权结果约为零点七五五、零点九一零。}}
\begin{tcolorbox}[enhanced,colback=paper,colframe=linegray,
boxrule=.7pt,arc=2mm,left=3mm,right=3mm,top=3mm,bottom=3mm,
coltitle=ink,fonttitle=\sffamily\bfseries,
title={三条 Key 的匹配分数与读取权重}]
\small
“开了”的 Query 为 $(1,1)$；Key 经独立投影，Value 仍是输入。
\begin{center}
\begin{tikzpicture}[x=1cm,y=.9cm,>=Latex]
\node[anchor=west,font=\small\bfseries] at (0,.5) {Token};
\node[anchor=west,font=\small\bfseries] at (2.6,.5) {点积分数};
\node[anchor=west,font=\small\bfseries] at (5.1,.5) {Softmax 权重};
\node[anchor=west,font=\small\bfseries] at (10.5,.5) {Value};
\foreach \yy/\word/\s/\weight/\len/\vv in {
 0/春天/1/0.090/.38/{(1,0)},
 -1.0/花/2/0.245/1.03/{(0,1)},
 -2.0/开了/3/0.665/2.79/{(1,1)}
}{
 \node[anchor=west,font=\small] at (.1,\yy) {\word};
 \node[font=\small] at (3.5,\yy) {$\s$};
 \fill[blueback] (5.1,\yy-.18) rectangle (8.1,\yy+.18);
 \fill[teal] (5.1,\yy-.18) rectangle ({5.1+\len},\yy+.18);
 \draw[linegray] (5.1,\yy-.18) rectangle (8.1,\yy+.18);
 \node[anchor=west,font=\small] at (8.25,\yy) {$\weight$};
 \node[anchor=west,font=\small] at (10.65,\yy) {$\vv$};
}
\node[anchor=west,font=\small] at (.1,-2.9)
 {$0.090(1,0)+0.245(0,1)+0.665(1,1)\approx(0.755,0.910)$};
\end{tikzpicture}
\end{center}
\end{tcolorbox}

\begin{center}\captionof{figure}{独立 Key 投影下的第三行 Query：匹配、归一化与 Value 读出}\label{fig:attention-weights}\end{center}
\tagstructend
